% ---- pg-534 （书页 522） ----

\noindent 11.36\quad 参考练习 11.17. 只用淡水生长和海水生长两个变量作为协变量, 对\hyperref[tab:11.2]{表 11.2} 中的鲑鱼数据作逻辑斯蒂回归. 对 $\alpha=0.05$, 检验模型的显著性和每个协变量的显著性. 用拟合的函数, 根据法则(11-77), 将\hyperref[tab:11.2]{表 11.2} 中的观测分为阿拉斯加鲑鱼和加拿大鲑鱼. 计算显示的误判率 APER, 并将此误判率与例 11.8 中得到的误判率作比较.

\phantomsection\label{bib:9}
\section*{参考文献}

\begin{enumerate}
\item Anderson, E. “The Irises of the Gaspé Peninsula.” \textit{Bulletin of the American Iris Society}, \textbf{59} (1939), 2--5.
\item Anderson, T. W. \textit{An Introduction to Multivariate Statistical Analysis} (3rd ed.). New York: John Wiley, 2003.
\item Bartlett, M. S. “An Inverse Matrix Adjustment Arising in Discriminant Analysis.” \textit{Annals of Mathematical Statistics}, \textbf{22} (1951), 107--111.
\item Bouma, B. N., et al. “Evaluation of the Detection Rate of Hemophilia Carriers.” \textit{Statistical Methods for Clinical Decision Making}, \textbf{7}, no. 2 (1975), 339--350.
\item Breiman, L., J. Friedman, R. Olshen, and C. Stone. \textit{Classification and Regression Trees}. Belmont, CA: Wadsworth, Inc., 1984.
\item Darroch, J. N., and J. E. Mosimann. “Canonical and Principal Components of Shape.” \textit{Biometrika}, \textbf{72}, no. 1 (1985), 241--252.
\item Efron, B. “The Efficiency of Logistic Regression Compared to Normal Discriminant Analysis.” \textit{Journal of the American Statistical Association}, \textbf{81} (1975), 321--327.
\item Eisenbeis, R. A. “Pitfalls in the Application of Discriminant Analysis in Business, Finance and Economics.” \textit{Journal of Finance}, \textbf{32}, no. 3 (1977), 875--900.
\item Fisher, R. A. “The Use of Multiple Measurements in Taxonomic Problems.” \textit{Annals of Eugenics}, \textbf{7} (1936), 179--188.
\item Fisher, R. A. “The Statistical Utilization of Multiple Measurements.” \textit{Annals of Eugenics}, \textbf{8} (1938), 376--386.
\item Ganesalingam, S. “Classification and Mixture Approaches to Clustering via Maximum Likelihood.” \textit{Applied Statistics}, \textbf{38}, no. 3 (1989), 455--466.
\item Geisser, S. “Discrimination, Allocatory and Separatory, Linear Aspects.” In \textit{Classification and Clustering}, edited by J. Van Ryzin, pp. 301--330. New York: Academic Press, 1977.
\item Gerrild, P. M., and R. J. Lantz. “Chemical Analysis of 75 Crude Oil Samples from Pliocene Sand Units, Elk Hills Oil Field, California.” \textit{U.S. Geological Survey Open-File Report}, 1969.
\item Gnanadesikan, R. \textit{Methods for Statistical Data Analysis of Multivariate Observations} (2nd ed.). New York: Wiley-Interscience, 1997.
\item Habbema, J. D. F., J. Hermans, and K. Van Den Broek. “A Stepwise Discriminant Analysis Program Using Density Estimation.” In \textit{Compstat 1974, Proc. Computational Statistics}, pp. 101--110. Vienna: Physica, 1974.
\item Hills, M. “Allocation Rules and Their Error Rates.” \textit{Journal of the Royal Statistical Society (B)}, \textbf{28} (1966), 1--31.
\item Hosmer, D. W. and S. Lemeshow. \textit{Applied Logistic Regression} (2nd ed.). New York: Wiley-Interscience, 2000.
\item Hudlet, R., and R. A. Johnson. “Linear Discrimination and Some Further Results on Best Lower Dimensional Representations.” In \textit{Classification and Clustering}, edited by J. Van Ryzin, pp. 371--394. New York: Academic Press, 1977.
\end{enumerate}
